\section{Discussion and Threats to Validity}

Internal validity threats include parser incompleteness, missing generated
binding outputs, toolchain/config differences, and oracle limitations. \tool{}
mitigates these threats with extraction diagnostics, per-version environment
metadata, config hashes, strict extractor audit, and explicit separation
between \primary{} and auxiliary build/wrapper evidence.

Construct validity threats include semantic label ambiguity and the fact that
warnings are review targets. The review guide separates \truesemantic{},
\truewrapper{}, benign drift, false positives, and unclear cases. The review
artifacts are treated as a trusted binddrift-review expert protocol because
the artifact records role separation, rank/score blindness, adjudication
coverage, agreement, and label-leakage checks. Overall warning-set precision is
low; the method targets prioritization, and the main evidence is Top-K yield,
baseline lift, false-positive taxonomy, and case-study evidence.

External validity threats include the Linux mainline scope, the x86\_64 main
labeled evaluation, and the Rust-for-Linux surface selection. The arm64 slice
reduces this threat by showing the replay and prioritization pipeline on a
second architecture, but it does not remove the threat or add independent
review labels.
